\section{Theta differentials and a nonzero mixed period}\label{sec:theta}

The surface criterion of Proposition~\ref{prop:surface-switch} asks for
a nonzero period made from four one-forms on one surface.  We construct
the analytic part of that input using theta series attached to auxiliary
Hermitian spaces of dimensions one and two, called multiplicity spaces.
For a multiplicity line, a component of the first Taylor tensor is a
holomorphic one-form.  Varying the finite theta input makes its values
span a cotangent fiber, so two such constructions give a nonzero wedge.
For their product input, the second Taylor tensor for the orthogonal
sum of the two lines has a double-alternating projection that is a
nonzero multiple of that wedge.  This total two-form allows us to change
the two multiplicity lines while keeping their sum and its finite theta
input fixed.

For the two sign pairs in the normalized switch, we will obtain total
two-forms \(\xi,\xi'\) on one fixed compact ball quotient \(X\).
On a common principal cover \(\pi:X'\to X\), they have finite
expansions into holomorphic one-forms:
\[
 \pi^*\xi=\sum_r a_r\Theta_1(\phi_{1,r})\wedge\Theta_2(\phi_{2,r}),
 \qquad
 \pi^*\xi'=\sum_s b_s\Theta_3(\phi_{3,s})\wedge\Theta_4(\phi_{4,s})
\]
such that
\begin{equation}\label{eq:theta-mixed-period-goal}
 \int_{X'} \pi^*\xi\wedge\overline{\pi^*\xi'}\ne0.
\end{equation}
The comparison of \(\xi'\) with \(\xi\) keeps the rational
symplectic space, the rational special-unitary group, the total finite
theta input, and the compact quotient fixed.  Only afterward will we
factor the rational endpoints and pass to \(X'\), on which the
individual one-forms in both sums descend.  At this stage they give
complex cohomology classes.  Section~\ref{sec:arithmetic} will provide
projective principal levels, and Section~\ref{sec:cm-sources} will
identify the rational CM source maps.

\subsection{The ball and the sign data}\label{subsec:theta-ball}

Let \(E\subset\Qbar\) be a Galois CM field, put \(F=E^+\), and assume
\(s=[F:\Q]\ge2\).  Write \(G_E=\Gal(E/\Q)\) and let \(c\in G_E\)
be complex conjugation.  Fix a set \(\Sigma\subset G_E\) of
representatives modulo \(c\), including the distinguished embedding
\(\mu=1\).  The construction for a single multiplicity line takes an
odd sign function \(v:G_E\to\{1,-1\}\) with \(v(1)=1\).  We index
its signs by
\begin{equation}\label{eq:theta-sign-inversion}
 d(\tau)=v(\tau^{-1}),\qquad
 d(c\tau)=-d(\tau),\qquad d(\mu)=1.
\end{equation}
Thus \(d(g^{-1})=v(g)\), the indexing needed by the arithmetic
calculation.  For the normalized switch, fix four such functions with
\(d_1+d_2=d_3+d_4\).  We first construct the one-form for an arbitrary
one of them, and then compare the pairs \((d_1,d_2)\) and
\((d_3,d_4)\).  All Hermitian forms in this section are linear in
their first argument.

Choose an \(E/F\)-Hermitian space \(V\) of \(E\)-dimension three,
of signature \((2,1)\) at \(\mu\) and positive definite at every
other real place of \(F\).  For example, take a diagonal form with two
totally positive coefficients and a third coefficient negative at
\(\mu\) and positive elsewhere; weak approximation in \(F\) supplies
the last coefficient.  Put
\[
 G=\SU(V),\qquad
 D=\{\,L\subset V_\mu:\dim_\C L=1,\ h_V|_L<0\,\}.
\]
The group \(G\) is defined over \(F\), and \(D\) is the complex
two-ball.  We also write \(L\) for the tautological negative line
bundle over \(D\), and set \(Q=\underline{V_\mu}/L\).  The
Grassmannian tangent calculation gives
\begin{equation}\label{eq:theta-cotangent}
 T_D=\Hom(L,Q),\qquad \Omega_D^1=L\otimes Q^*.
\end{equation}

Fix an \(\mathcal O_E\)-lattice \(\Lambda\subset V\); all principal congruence
levels below are matrix levels for this lattice.  Their quotients of
\(D\) are compact.  To see the relevant anisotropy, an isotropic
vector defined over \(E\) would still be isotropic at a definite real
completion, which is impossible.  Equivalently, \(G\) has no
nontrivial \(F\)-split torus: a nonzero weight space of such a torus
in \(V\) would be isotropic, since invariance of \(h_V\) forces
\(h_V(x,x)=0\) for a vector of nontrivial split weight.  Thus \(G\)
is \(F\)-anisotropic.  The arithmetic lattice and compactness
theorems therefore make its arithmetic subgroups cocompact
lattices in \(G(F\otimes_\Q\R)\)
\cite{BorelHarishChandra1962}; see also the formulation in
\cite[Theorem~3.3 and Proposition~3.5]{MilneShimura2017}.  Here
\[
 G(F\otimes_\Q\R)\simeq \SU(2,1)\times\prod_{\tau\ne\mu}\SU(3).
\]
Projection past the compact factors preserves discreteness and
cocompactness.  Indeed, a sequence in an arithmetic subgroup whose
active projections approach the identity has, after taking a
subsequence, convergent compact components; differences would then
approach the identity in the full discrete subgroup.  A sufficiently
deep principal matrix level is torsion-free.  Its action on \(D\) is
free, because a discrete subgroup has finite intersection with a
point stabilizer.  Hence its quotient is a smooth connected compact
K\"ahler surface.  Later we will impose a principal level for which
the reference moduli construction also proves projectivity.

\subsection{Positive Lagrangians and differential bundles}\label{subsec:theta-lagrangians}

Fix \(\delta\in E^\times\) with \(c(\delta)=-\delta\).
For a nondegenerate \(E/F\)-Hermitian space \(P\) of \(E\)-dimension
\(n\in\{1,2\}\), let
\[
 \mathbb V_P=\Res_{E/F}(V\otimes_E P).
\]
It is an \(F\)-space of dimension \(6n\).  The product of the two
Hermitian forms is a well-defined Hermitian form on \(V\otimes_E P\),
and we use the \(F\)-valued form
\begin{equation}\label{eq:theta-symplectic-form}
 \psi_P(v\otimes p,v'\otimes p')
  =\Tr_{E/F}\!\bigl(\delta\,h_V(v,v')h_P(p,p')\bigr).
\end{equation}
Interchanging the arguments conjugates the Hermitian product and
changes the sign of \(\delta\), so \(\psi_P\) is alternating.  It
is nondegenerate: the tensor-product Hermitian form is nondegenerate,
and the trace pairing on \(E/F\) is nondegenerate.  More explicitly,
if the Hermitian pairing of \(x\ne0\) with some \(y\) is \(b\ne0\),
then its pairing with \(e y\), for \(e\in E\), is \(c(e)b\);
trace duality supplies \(e\) for which
\(\Tr_{E/F}(\delta c(e)b)\ne0\).  Thus \((\mathbb V_P,\psi_P)\) is
a symplectic space over \(F\).  Its underlying rational vector space
carries the alternating form \(\Tr_{F/\Q}\psi_P\); this is the
restriction of scalars used in complexifications over \(\Q\) below.
Symplectic coordinates and transformations will still be taken over
\(F\).

First take \(P=W_i\) to be a Hermitian line associated with a sign
function \(d_i\) as in \eqref{eq:theta-sign-inversion}.  Complexifying
the trace pairing identifies the \(c\tau\) summand complex-linearly
with the dual of the \(\tau\) summand.  Define a subbundle
\(\mathcal F_i\subset\mathbb V_P\otimes_\Q\C\) over \(D\) at the
active pair of embeddings by
\begin{equation}\label{eq:theta-active-lagrangian}
 (\mathcal F_i)_\mu=L\otimes(W_i)_\mu,\qquad
 (\mathcal F_i)_{c\mu}=Q^*\otimes(W_i)_\mu^*.
\end{equation}
Here \(Q^*=\operatorname{ann}(L)\subset V_\mu^*\), embedded in
the conjugate summand by the Hermitian pairing.  It is the
holomorphic annihilator bundle, not an antiholomorphically varying
orthogonal complement in \(V_\mu\).  At
\(\tau\in\Sigma\setminus\{\mu\}\), take the full \(c\tau\) summand
when \(d_i(\tau)=1\), and the full \(\tau\) summand when
\(d_i(\tau)=-1\); take zero on the other summand.  For a rational
orthogonal sum \(P=W_1\perp W_2\), take the direct sum of these
prescriptions and write \(\mathcal F_P=\mathcal F_1\oplus\mathcal F_2\).

\begin{lemma}\label{lem:theta-positive-lagrangian}
The Hermitian forms on the lines \(W_i\) can be chosen over \(E/F\)
so that \(\mathcal F_P\) is a holomorphic positive Lagrangian for
\(\psi_P\), equivariantly for \(G(F\otimes_\Q\R)\).  Its determinant
and an equivariant square root have the form
\begin{equation}\label{eq:theta-halfdet}
 \det\mathcal F_P\simeq L^{2n}\otimes C_P,\qquad
 (\det\mathcal F_P)^{1/2}\simeq L^n\otimes C_P^{1/2},
\end{equation}
where \(C_P\) is a constant complex line on which
\(G(F\otimes_\Q\R)\) acts trivially.
\end{lemma}

\begin{proof}
Write \(\tau(\delta)=it_\tau\), where \(t_\tau\in\R^\times\).
At the real place represented by \(\tau\in\Sigma\), write
\(\psi_{P,\tau}\) for the complex bilinear extension of the
real-place form.  Our positivity convention is
\[
 \psi_{P,\tau}(u,\overline u)/i>0
 \qquad\bigl(0\ne u\in
       (\mathcal F_P)_\tau\oplus(\mathcal F_P)_{c\tau}\bigr).
\]
On the \(\tau\) summand its sign is the sign of
\(t_\tau h_Vh_P\); on the \(c\tau\) summand the sign is the
opposite.  At \(\mu\), choose the coefficient of each line form to
have sign \(-\operatorname{sgn}(t_\mu)\).  The negative line \(L\)
on the \(\mu\) side and the dual of the positive plane on the
\(c\mu\) side then both have positive sign.  At a definite place
\(\tau\ne\mu\), choose the line coefficient to have sign
\[
 -d_i(\tau)\operatorname{sgn}(t_\tau).
\]
The chosen full summand has positive sign in either case.
These finitely many strict sign conditions on an element of \(F\)
are simultaneously realizable by weak approximation.

At the active pair the two subspaces are annihilators under the
complexified pairing; at every other pair one side is full and the
other zero.  Thus \(\mathcal F_P\) is isotropic of half the total
dimension.  The active annihilator description proves holomorphy,
and the compact-place summands are constant over \(D\).
Equivariance follows from the action on \(L\) and its annihilator
at \(\mu\), and from preservation of the full embedding summands
at the compact places.

Finally
\[
 \det Q^*=L\otimes\det(V_\mu)^*.
\]
For general \(P\) the active determinant is therefore
\[
 \det(L\otimes P_\mu)\otimes\det(Q^*\otimes P_\mu^*)
 \simeq
 L^{2n}\otimes\bigl(\det(V_\mu)^*\bigr)^{\otimes n}
       \otimes\det(P_\mu)^*.
\]
All factors other than \(L^{2n}\), including the determinants at
the compact places, are constant lines.  Their special-unitary
actions are trivial, because they involve powers of determinants
of \(V_\tau\) or \(V_\tau^*\), while \(G\) acts trivially on
\(P\).  This proves the first isomorphism in
\eqref{eq:theta-halfdet}; choosing a square root of the constant
line gives the second.
\end{proof}

We choose nonzero complex vectors in these constant lines when
identifying tensors with differential forms.  Such choices only
rescale those forms by nonzero constants.  No rational trivialization
is asserted.  Also, the constant lines are trivial for \(G=\SU(V)\);
they need not be trivial under the larger unitary group.

The determinant calculation gives the differential targets for the
theta tensors.  For a multiplicity line \(W_i\), the relevant target is
\begin{equation}\label{eq:theta-one-target}
\begin{aligned}
 (\det\mathcal F_i)^{1/2}\otimes(\mathcal F_i)_{c\mu}
 &\simeq L\otimes Q^*\otimes(\text{constant line})\\
 &\simeq \Omega_D^1\otimes(\text{constant line}).
\end{aligned}
\end{equation}
For \(W=W_1\perp W_2\), define the double-alternating map
\[
 \operatorname{alt}_2:
 \Sym^2(Q^*\otimes W_\mu^*)\longrightarrow
       \Lambda^2Q^*\otimes\Lambda^2W_\mu^*
\]
by
\[
 \operatorname{alt}_2\bigl((a\otimes u)(b\otimes v)\bigr)
       =(a\wedge b)\otimes(u\wedge v).
\]
Here the product on the left is multiplication in the symmetric
algebra.  Exchanging its two factors changes both exterior signs,
so the map is well-defined.  Including the half-determinant line gives
\begin{equation}\label{eq:theta-top-target}
 (\det\mathcal F_W)^{1/2}\otimes
       \Lambda^2Q^*\otimes\Lambda^2W_\mu^*
       \simeq K_D\otimes(\text{constant line}),
\end{equation}
because \(K_D=\det\Omega_D^1=L^2\otimes\Lambda^2Q^*\).
The first theta tensor will give the one-form in
\eqref{eq:theta-one-target}; the double-alternating part of the second
will give the two-form in \eqref{eq:theta-top-target}.

The reason to combine the two lines is geometric.  At the active pair,
their total positive Lagrangian has the pieces
\[
 L\otimes W_\mu,\qquad Q^*\otimes W_\mu^*.
\]
These depend only on \(W_\mu\), independently of its decomposition
into lines.  At a compact place, let \(W_\tau^+\) and \(W_\tau^-\)
be the sums of the multiplicity lines with respective signs \(+1\)
and \(-1\) under \(d_i\).  The total positive pieces there are
\[
 V_\tau\otimes W_\tau^-,\qquad
 V_\tau^*\otimes(W_\tau^+)^*;
\]
the second dual is embedded as
\(\operatorname{ann}(W_\tau^-)\subset W_\tau^*\).
The relation \(d_1+d_2=d_3+d_4\) says that the two sign pairs have
the same multiset at each place.  Thus the original local lines can
be relabeled with signs \(d_3,d_4\) without changing these total
pieces.  Weak approximation will produce a rational decomposition
with the new signs and nearby compact-place data, while leaving the
active total Lagrangian unchanged.  We will apply this approximation
in Subsection~\ref{subsec:theta-approximation} to one fixed theta
input.  To compare the resulting forms on a compact quotient, we
first establish their covariance and congruence descent.

\subsection{Rational theta covariance}\label{subsec:theta-covariance}

Choose symplectic coordinates over \(F\) on \(\mathbb V_P\),
written \((p,q)\in F^a\oplus F^a\), where \(a=3n\), with
\[
 \psi_P((p,q),(p',q'))=p^tq'-q^tp'.
\]
The Schr\"odinger position coordinates are \(q\).  A positive
Lagrangian at a real place has a unique graph frame
\[
 e_Z=\binom ZI,\qquad Z=Z^t,\qquad \operatorname{Im}Z>0.
\]
Indeed, it cannot meet the \(p\)-axis nontrivially, because
\(\psi_P(u,\bar u)\) vanishes there.  Projection to the
\(q\)-axis is therefore an isomorphism.  Isotropy says \(Z^t=Z\),
and, for a graph vector,
\[
 \psi_{P,\tau}(e_Zu,\overline{e_Zu})/i
       =2u^t(\operatorname{Im}Z)\bar u,
\]
which proves the last condition.  The positive Lagrangians above
thus give a period map \(D\to\prod_{\tau\in\Sigma}\mathfrak H_a\).

Let \(\phi\in\mathcal S(\A_{F,f}^a)\) be a finite-adelic
Schwartz--Bruhat function.  For
\(Z=(Z_\tau)_\tau\in\prod_\tau\mathfrak H_a\) and
\(z=(z_\tau)_\tau\in\bigoplus_\tau\C^a\), put
\begin{equation}\label{eq:theta-series}
 \vartheta_\phi(Z,z)
  =\sum_{x\in F^a}\phi(x)
    \exp\!\left(
       \pi i\sum_{\tau\in\Sigma}x_\tau^tZ_\tau x_\tau+
       2\pi i\sum_{\tau\in\Sigma}x_\tau^tz_\tau\right).
\end{equation}
The rational transformation law for this series includes an exponential
quadratic in \(z\), which mixes its ordinary Taylor coefficients.
The following completion cancels that factor and gives coefficients
with a tensor transformation law.  Define
\begin{equation}\label{eq:theta-completion}
 \begin{split}
 \widehat\vartheta_\phi(Z,z)
  &=\exp\!\left(\frac{\pi}{2}\sum_{\tau\in\Sigma}
       z_\tau^t(\operatorname{Im}Z_\tau)^{-1}z_\tau\right)
       \vartheta_\phi(Z,z),\\
 \widehat\vartheta_\phi(Z,z)
  &=\sum_{k\ge0}\langle T_{k,\phi}(Z),z^{\otimes k}\rangle .
 \end{split}
\end{equation}
Thus \(T_{k,\phi}\) includes the factor \(1/k!\) in the \(k\)-th
derivative.  We use multiplication in the symmetric algebra for
these homogeneous coefficients.

The series and every fixed order of its derivatives converge
uniformly on compact subsets of the Siegel and \(z\) domains.
Indeed, the rational points in the support of \(\phi\) lie in
finitely many translates of a full lattice in
\((F\otimes_\Q\R)^a\).  On a compact set the imaginary parts have
a common positive lower bound, and the summands, after a fixed
number of differentiations, are bounded by a polynomial in
\(\|x\|\) times
\[
 C\exp(-c\|x\|^2+C'\|x\|).
\]
The lattice sum of this bound converges.  The completion and its
derivatives obey the same bounds on compact sets.  The same
conclusion holds after any fixed number of derivatives in smooth
parameters on a compact subset of a smooth parameter domain:
there the period map and all derivatives of it needed for that
order are bounded, and its imaginary parts have a common positive
lower bound.  Compactness of the image of a period map alone,
without bounds on its parameter derivatives, would not imply this
last assertion.

We next specify the rational Weil-representation input, including
the cover being used.  The circle extension is obtained from the
metaplectic double cover by extending its central \(\{1,-1\}\)
to the group of unit complex scalars.  Take the global additive character of
\(\A_F/F\) whose infinite component is
\(t\mapsto\exp(2\pi i\Tr_{F/\Q}t)\).  Weil's rational section in
the adelic circle extension leaves the summation functional
\[
 \lambda(\Phi)=\sum_{x\in F^a}\Phi(x)
\]
invariant.  In characteristic zero his pseudosymplectic group is
the symplectic group, and his adelic twofold subgroup contains this
rational section
\cite[Ch.~III, nos.~31 and 41, Eq.~(39) and Theorem~6;
      Ch.~IV, no.~45, pp.~200--201]{Weil1964}.
Consequently, for a rational symplectic transformation we may
choose an infinite double-cover lift and its paired finite lift
whose adelic product is the rational lift.  A change of the
infinite sign requires the same change of the finite sign.
These paired lifts give rational covariance; they are not an
assertion that the symplectic double cover splits on an arbitrary
adelic unitary group.

\begin{lemma}\label{lem:theta-tensor-law}
Let \(g\in\Sp(\mathbb V_P)(F)\).  At each real place write
\[
 g=\begin{pmatrix}A&B\\C&D_0\end{pmatrix},\qquad
 J=CZ+D_0,\qquad gZ=(AZ+B)J^{-1}.
\]
Choose an infinite metaplectic lift with factor \(j(g,Z)\) satisfying
\(j(g,Z)^2=\prod_{\tau\in\Sigma}\det J_\tau\), and let
\(\widetilde g_f\) be its paired finite lift.  Then
\begin{equation}\label{eq:theta-jacobi-law}
 \vartheta_{\omega_f(\widetilde g_f)\phi}(gZ,J^{-t}z)
  =j(g,Z)\exp\!\left(\pi i\sum_{\tau\in\Sigma}
             z_\tau^tJ_\tau^{-1}C_\tau z_\tau\right)
             \vartheta_\phi(Z,z).
\end{equation}
The completed coefficients satisfy
\begin{equation}\label{eq:theta-taylor-law}
 T_{k,\omega_f(\widetilde g_f)\phi}(gZ)
   =j(g,Z)\Sym^k(J)\,T_{k,\phi}(Z).
\end{equation}
They are thus tensors valued in
\[
 (\det\mathcal F_P)^{1/2}\otimes\Sym^k\mathcal F_P.
\]
In particular the positive Lagrangian, not its dual, occurs in
the symmetric factor.  The tensor \(T_{1,\phi}\) is holomorphic.
\end{lemma}

\begin{proof}
We first check the infinite formula, keeping track of its signs.
At one real place write
\[
 f_{Z,z}(x)=\exp(\pi ix^tZx+2\pi ix^tz),\qquad K(g,Z)=J^{-1}C.
\]
The matrix \(K(g,Z)\) is symmetric; this follows from
\(JC^t=CJ^t\), a symplectic identity.  The formula to prove is
\begin{equation}\label{eq:theta-gaussian-law}
 \omega_\infty(\widetilde g)f_{Z,z}
  =j(g,Z)^{-1}\exp(-\pi iz^tK(g,Z)z)
       f_{gZ,J^{-t}z}.
\end{equation}
For symmetric \(B\), the upper unipotent
\(n(B)=\left(\begin{smallmatrix}I&B\\0&I\end{smallmatrix}\right)\)
acts by multiplication by \(\exp(\pi ix^tBx)\), which gives
\eqref{eq:theta-gaussian-law} with \(j=1\).  For
\(m(A)=\diag(A,A^{-t})\), the chosen lift acts by
\[
 \omega_\infty(\widetilde m(A))b(x)
       =j(m(A))^{-1}b(A^tx),\qquad
 j(m(A))^2=(\det A)^{-1}.
\]
For \(\det A>0\) the positive choice has
\(j(m(A))^{-1}=|\det A|^{1/2}\).  For negative determinant
the metaplectic phase is retained in \(j\).  Since
\(m(A)Z=AZA^t\) and \(J^{-t}z=Az\), this is the required formula.

For
\(w=\left(\begin{smallmatrix}0&-I\\I&0\end{smallmatrix}\right)\),
choose the lift
\[
 b\longmapsto e^{-\pi ia/4}
        \int_{\R^a}b(y)e^{-2\pi ix^ty}\,dy.
\]
The Gaussian integral is
\[
 \int_{\R^a}f_{Z,z}(y)e^{-2\pi ix^ty}\,dy
  =\det(-iZ)^{-1/2}
       \exp\!\bigl(-\pi i(x-z)^tZ^{-1}(x-z)\bigr),
\]
where the square root of \(\det(-iZ)\) equals \(1\) at \(Z=iI\).
Taking
\[
 j(w,Z)=e^{\pi ia/4}\det(-iZ)^{1/2}
\]
gives \(j(w,Z)^2=\det Z\) and proves
\eqref{eq:theta-gaussian-law} for \(w\).

These generators give the formula for the whole real symplectic
group.  To check that the displayed exponent also composes
correctly, let \(J_2=J(g_2,Z)\) and
\(J_1=J(g_1,g_2Z)\).  The symplectic identities give
\[
 J(g_1g_2,Z)=J_1J_2,\qquad
 K(g_1g_2,Z)=K(g_2,Z)+J_2^{-1}K(g_1,g_2Z)J_2^{-t}.
\]
For the second identity, use
\(A_2-(g_2Z)C_2=J_2^{-t}\) to write
\(C_{12}=C_1J_2^{-t}+J_1C_2\), and multiply by
\((J_1J_2)^{-1}\).  Together with the metaplectic group law for
\(j\), this is exactly the composition law in
\eqref{eq:theta-gaussian-law}.  Multiplying the formula over the
real places gives its infinite adelic version.

Apply rational invariance of \(\lambda\) to
\(\phi\otimes\prod_\tau f_{Z_\tau,z_\tau}\).  By
\eqref{eq:theta-gaussian-law}, the transformed sum is
\[
 j(g,Z)^{-1}
 \exp\!\left(-\pi i\sum_\tau z_\tau^tJ_\tau^{-1}C_\tau z_\tau\right)
 \vartheta_{\omega_f(\widetilde g_f)\phi}(gZ,J^{-t}z).
\]
Equating this with the original sum gives
\eqref{eq:theta-jacobi-law}.  In particular the finite action on
the left is \(\omega_f(\widetilde g_f)\), with no inverse.
If one instead keeps the finite input on the left fixed, the input
on the right is acted on by \(\omega_f(\widetilde g_f)^{-1}\).

For the completion, put \(Y=\operatorname{Im}Z\) and
\(Y'=\operatorname{Im}(gZ)\) at one place.  From
\(Y'=\bar J^{-t}YJ^{-1}\), and then
\(\bar J^t=J^t-2iYC^t\), we obtain
\begin{equation}\label{eq:theta-completion-identity}
 J^{-1}(Y')^{-1}J^{-t}=Y^{-1}-2iJ^{-1}C.
\end{equation}
Here \(C^tJ^{-t}=J^{-1}C\) by symmetry.  Thus the completion
factor at \((gZ,J^{-t}z)\) differs from that at \((Z,z)\) by
\(\exp(-\pi iz^tJ^{-1}Cz)\).  This cancels the exponential in
\eqref{eq:theta-jacobi-law}, so
\[
 \widehat\vartheta_{\omega_f(\widetilde g_f)\phi}(gZ,J^{-t}z)
       =j(g,Z)\widehat\vartheta_\phi(Z,z).
\]
Comparison of homogeneous terms yields
\eqref{eq:theta-taylor-law}.  Finally,
\[
 g e_Z=e_{gZ}J.
\]
Hence \(J\) is the transformation matrix for vectors in the
positive Lagrangian; this proves the asserted tensor target and
its orientation.  The completion has no linear term in \(z\), so
\(T_{1,\phi}\) is the ordinary holomorphic theta gradient.
For comparison, the first-order law agrees with the classical
theta-gradient formula in \cite[\S9.1]{MollerZagier2016}.
\end{proof}

For a multiplicity line \(W_i\), the \(c\mu\)-projection of
\(T_{1,\phi}\) therefore lies in the cotangent target
\eqref{eq:theta-one-target}.  Rational covariance alone does not
supply a congruence subgroup fixing a given finite input.
Subsection~\ref{subsec:theta-splittings} obtains such stabilizers
from an honest smooth finite-adelic action.

\begin{remark}\label{rem:theta-adapted-coordinates}
We will use the lemma to change rational Schr\"odinger coordinates.
Suppose a rational symplectic matrix \(S\) takes adapted coordinate
vectors to the reference coordinate vectors.  Then
\(Z_{\mathrm{ref}}=S Z_{\mathrm{ad}}\), and the input in the
adapted coordinates is
\[
 \phi_{\mathrm{ad}}
    =\omega_f(\widetilde S_f)^{-1}\phi_{\mathrm{ref}}.
\]
Equation~\eqref{eq:theta-taylor-law} compares the resulting tensors.
Changing the sign of the paired lift changes both the infinite
factor and the finite input sign, so their comparison is unchanged.
For coordinates adapted to a rational symplectic direct sum, the
period matrix is block diagonal and the completed theta series of
a product input is the product of the two completed series.
\end{remark}

\subsection{Unitary splittings and congruence descent}\label{subsec:theta-splittings}

The preceding lemma uses paired lifts of rational symplectic
elements.  An honest action on finite inputs, and therefore a
congruence stabilizer for each input, requires a different statement.
We use the unitary splitting theorem in the following precise form.
For a quadratic extension \(E/F\), a Hermitian space \(H\) of
dimension \(m\), a skew-Hermitian space \(K\) of dimension \(n\),
and a global additive character, choose unitary Hecke characters
\(\chi_H,\chi_K\) of \(E\) with
\[
 \chi_H|_{\A_F^\times/F^\times}=\epsilon_{E/F}^{\,m},\qquad
 \chi_K|_{\A_F^\times/F^\times}=\epsilon_{E/F}^{\,n}.
\]
The corresponding unitary actions lift to the \emph{circle-scalar
extension} of the adelic symplectic Weil representation.  The
global lift is assembled from local lifts and agrees with the
canonical rational symplectic section on rational unitary points.
Its finite part acts smoothly on finite Schwartz functions.
Changing the splitting characters twists the unitary actions by
characters composed with their determinants.  Here
\(\epsilon_{E/F}\) is the quadratic idele class character.

The local splitting is due to Kudla \cite{Kudla1994}.  We use the
circle-extension formulation and its local and adelic compatibility
in \cite[\S\S3.1--3.1.1 and 4.1--4.1.1]{BoradeEtAl2025ToriPeriods}.
In the present CM setting, adelic automorphy and the determinant-twist
rule are also recorded in \cite[\S7.2 and Remark~7.3]{BMM2016}.
Throughout this invocation the extension is the circle pushout:
we do not assert that the double cover splits on every unitary
group.  We will compare this circle lift with an actual double
lift only after restriction to \(G=\SU(V)\).
The split-place normalization below will also enter the local Hecke
calculation of Subsection~\ref{subsec:active-parameter}.

\begin{lemma}\label{lem:unitary-splitting}
For the symplectic space \((\mathbb V_P,\psi_P)\), restriction of
the unitary splitting to \(G(\A_F)\) yields an honest finite-adelic
Weil action \(\omega_{P,f}\).  It is compatible with the rational
covariance in Lemma~\ref{lem:theta-tensor-law} and is independent
of the splitting characters.  At infinity its half-determinant
action is the natural action on \(L^n\) in
\eqref{eq:theta-halfdet}.

At a split finite place, choose as Schr\"odinger positions for a
multiplicity line one standard three-dimensional embedding
summand.  The resulting \(\SL_3\)-action is inverse linear
substitution.  More generally, the normalized linear action
associated with a transformation \(h\) of the position space is
\begin{equation}\label{eq:theta-linear-action}
 (\rho(h)b)(x)=|\det h|^{-1/2}b(h^{-1}x).
\end{equation}
A bare metaplectic lift outside \(\SL_3\) is the normalized
operator \(\rho(h)\) times its lift scalar; that scalar must be
retained in any use of such a lift.

Every finite list of inputs is fixed under \(\omega_{P,f}\) by
all sufficiently deep principal lattice congruence subgroups.
\end{lemma}

\begin{proof}
First we check that the required splitting characters exist.
For the idele class groups \(C_F\) and \(C_E\), the natural map
\(C_F\to C_E\) is injective: if an \(F\)-idele becomes a principal
\(E\)-idele, its generator is fixed by \(c\) and hence belongs to
\(F\).  This map is also proper.  The absolute idele norm on
\(C_E\) restricts to the square of the norm on \(C_F\), whose
norm-one kernel is compact.  A compact subset of \(C_E\) has
norms in a compact subinterval of \(\R_{>0}\); its preimage is
therefore a closed subset of a compact norm interval times the
compact norm-one kernel of \(C_F\).  Thus \(C_F\) is a closed
locally compact subgroup of \(C_E\).  The character-extension
theorem for a closed subgroup of a locally compact abelian group
extends each unitary character \(\epsilon_{E/F}^{\,r}\) from
\(C_F\) to \(C_E\)
\cite[\S4.3, Corollary~4.42]{Folland2015}.  Apply this with
\(r=m,n\).

The splitting theorem applies to our form by replacing \(h_P\)
with the skew-Hermitian form \(\delta h_P\).  With both forms linear
in the first variable, the trace product in
\cite[\S\S2.1 and 3.1]{BoradeEtAl2025ToriPeriods} is exactly
\eqref{eq:theta-symplectic-form}.  We take the same global additive
character as in Lemma~\ref{lem:theta-tensor-law}.  The determinant twists from
changing the characters vanish on \(G=\SU(V)\).  The product of
the finite local splittings is therefore an honest action
\(\omega_{P,f}\) on \(G(\A_{F,f})\), independent of those
characters.

We compare its infinite circle lift with the double lift needed
for the differential bundles.  At the active maximal compact
subgroup \(S(\U(2)\times\U(1))\), let \(\ell\) be the character
on the negative line.  The determinant of the positive plane is
\(\ell^{-1}\), so the determinant character on
\(\mathcal F_P\) is \(\ell^{2n}\).  At every other real place
the determinant character is trivial on \(\SU(3)\).  On the
maximal compact \(\U(a)\) of each real symplectic factor, the
metaplectic double cover is the determinant-square-root cover.
For the product of the real factors its square is the product of
those determinant characters.  Its pullback to the maximal
compact subgroup of \(G(F\otimes_\Q\R)\) therefore splits with
square-root character \(\ell^n\).

A connected real reductive group retracts onto its maximal compact
subgroup.  The preceding determinant computation says that every
fundamental-group loop of \(G(F\otimes_\Q\R)\) has even
determinant winding in the symplectic group.  The lift from the
universal covering group consequently descends to a group lift
of \(G(F\otimes_\Q\R)\) into the double cover.  Its
half-determinant action is the one on \(L^n\) in
\eqref{eq:theta-halfdet}.  The original circle lift and this
double lift, pushed out to the circle extension, differ by a
continuous character of \(G(F\otimes_\Q\R)\).  This group is
connected and semisimple.  The differential of such a character
is zero, and connectedness makes the character trivial.
The two lifts therefore agree.  The rational compatibility of
the global splitting now pairs \(\omega_{P,f}(\gamma)\) with
this infinite lift for every \(\gamma\in G(F)\), proving the
claimed rational covariance.

At a split finite place the two embedding summands are dual
stable Lagrangians.  If \(h\) is the action on the position
summand, the symplectic matrix is \(\diag(h^{-t},h)\).
The Schr\"odinger linear formula gives
\eqref{eq:theta-linear-action}, up to the scalar of the chosen
bare lift.  On \(\SL_3(F_v)\), the normalized substitution lift
and the unitary lift have a character as their ratio.  This
character is trivial because \(\SL_3(F_v)\) is perfect:
elementary matrices generate it, and for distinct \(i,j,k\),
\[
 e_{ij}(u)=[e_{ik}(u),e_{kj}(1)].
\]
The determinant factor is also \(1\) there.  This proves the
split-place assertion without removing scalars from lifts outside
\(\SL_3\).

Finally the finite Weil action is smooth, so finitely many
Schwartz functions have a common open stabilizer.  Intersect it
with the fixed lattice stabilizer in \(G(\A_{F,f})\).
Principal lattice congruence neighborhoods form a basis at the
identity in that stabilizer, so a sufficiently deep principal
one lies in the intersection.  This proves the last assertion.
\end{proof}

Restrict the tensors of Lemma~\ref{lem:theta-tensor-law} to the
period map from \(D\).  If a finite list of inputs is fixed by an
open compact \(K_f\subset G(\A_{F,f})\), the lemma just proved
and \eqref{eq:theta-halfdet} give their bundle covariance for
\(\Gamma=G(F)\cap K_f\).  Passing to a torsion-free principal
subgroup makes all of them descend to one compact quotient.
This uses stabilizers in the smooth finite-adelic representation;
it makes no claim that an arbitrary finite-index subgroup of a
ball lattice is congruence.

\subsection{Theta one-forms span cotangent fibers}\label{subsec:theta-one-forms}

For a line \(W_i\), project \(T_{1,\phi}\) to the \(c\mu\)
summand of the positive Lagrangian.  Using our fixed complex
trivialization of the constant line in \eqref{eq:theta-one-target},
denote the resulting one-form on
\(D\) by \(\Theta_i(\phi)\).

\begin{proposition}\label{prop:theta-one}
For every finite input \(\phi\), the form \(\Theta_i(\phi)\) is
holomorphic on \(D\) and descends to every sufficiently deep
principal congruence quotient.  At any fixed point of \(D\), the
values \(\Theta_i(\phi)\), as \(\phi\) varies over finite
Schwartz inputs, span the cotangent fiber.  Finitely many actual
inputs already span that fiber and descend at a common sufficiently
deep principal level.
\end{proposition}

\begin{proof}
Holomorphy follows from that of \(T_{1,\phi}\), the holomorphic
embedding projection, and \eqref{eq:theta-cotangent}.  Descent
follows from Lemma~\ref{lem:unitary-splitting}.

Fix a point \(x_*\in D\) and its period point \(Z\).  Here \(a=3\).  For
\(0\ne x_0\in\mathcal O_F^a\) and positive integers \(k\), take
\[
 \phi_k=\mathbf1_{\,x_0+k\widehat{\mathcal O}_F^{\,a}}.
\]
Its rational support is exactly \(x_0+k\mathcal O_F^a\).
The first Taylor tensor is the ordinary gradient, so in the
chosen graph frame
\[
 T_{1,\phi_k}(Z)
  =2\pi i\sum_{x\in x_0+k\mathcal O_F^a}
      (x_\tau)_{\tau\in\Sigma}
       \exp\!\left(\pi i\sum_\tau x_\tau^tZ_\tau x_\tau\right).
\]
Under the Minkowski embedding \(\mathcal O_F^a\) is a full
lattice.  If \(\rho>0\) is the minimum norm of a nonzero lattice
vector, then, for \(m\ne0\) in this lattice and
\(k\ge2\|x_0\|/\rho\),
\[
 \|x_0+km\|\ge k\|m\|/2.
\]
The Gaussian bound used above, with any polynomial factor
absorbed into a weaker Gaussian, bounds the tail by a constant
times a convergent lattice sum of the form
\(\sum_{m\ne0}\exp(-c k^2\|m\|^2)\).  That sum tends to zero.
The same argument is uniform for \(Z\) in a compact
neighborhood.  We have proved
\begin{equation}\label{eq:theta-sparse-limit}
 T_{1,\phi_k}(Z)\longrightarrow
  2\pi i\,(x_{0,\tau})_{\tau\in\Sigma}
       \exp\!\left(\pi i\sum_\tau
                       x_{0,\tau}^tZ_\tau x_{0,\tau}\right).
\end{equation}

The scalar exponential is nonzero.  As \(x_0\) varies, the
vectors \((x_{0,\tau})_\tau\) span the real position space,
and their complex span is the entire first-tensor fiber in this
frame.  The complex span of the actual values \(T_{1,\phi}(Z)\)
is a subspace of a finite-dimensional space and hence is closed.
It contains all the limits in \eqref{eq:theta-sparse-limit},
so it is the full fiber.  Projection onto
\((\det\mathcal F_i)^{1/2}\otimes(\mathcal F_i)_{c\mu}\)
is surjective; therefore the projected values span
\(\Omega_D^1|_{x_*}\).  Choose finitely many actual inputs whose
values contain a basis, and then apply the stabilizer assertion
of Lemma~\ref{lem:unitary-splitting}.

The inputs \(\phi_k\) in this limit need not have a common
congruence stabilizer as \(k\) varies.  The limit takes place
in one fiber over \(D\); a common level is chosen only after
finitely many actual inputs have been selected.
\end{proof}

On a compact K\"ahler quotient a holomorphic one-form is closed
and harmonic.  A nonzero such form therefore has nonzero
cohomology class.  These are complex classes.  Nothing in this
proposition identifies a rational CM source for them.

\subsection{Rank-two tensors, wedges, and fixed-data variation}\label{subsec:theta-rank-two}

Take \(W=W_1\perp W_2\), so \(n=2\).  Project \(T_{2,\phi}\)
to \(\Sym^2((\mathcal F_W)_{c\mu})
=\Sym^2(Q^*\otimes W_\mu^*)\), and then apply
\(\operatorname{alt}_2\).  Its target is the canonical
bundle in \eqref{eq:theta-top-target}, up to a constant line.
After trivializing that line, denote the resulting differential
two-form by \(\Xi_W(\phi)\).  At a rational decomposition the product
formula for theta series will express it as a finite sum of wedges
of the one-forms just constructed.  We then vary the compact-place
data of the total space to compare two such expressions.

We now specify exactly what may vary in this construction.
Keep fixed the rational Hermitian space \(W\), the symplectic
space \((\mathbb V_W,\psi_W)\), the rational embedding
\(G\to\Sp(\mathbb V_W)\), reference rational Schr\"odinger
coordinates, the global Weil and unitary splitting choices, and
the resulting finite action \(\omega_{W,f}\).  Also fix a single
total finite input \(\phi\) in these reference coordinates.
At a compact place \(\tau\ne\mu\), let
\[
 W_\tau=W_\tau^+\perp W_\tau^-
\]
range over orthogonal decompositions with
\[
 \dim W_\tau^\pm=\#\{i\in\{1,2\}:d_i(\tau)=\pm1\}.
\]
Require the form on \(W_\tau^+\) to have sign
\(-\operatorname{sgn}(t_\tau)\), and the form on
\(W_\tau^-\) the opposite sign.  Let \(\mathcal T\) be the
product of these parameter spaces.  A definite local factor is
a point; when the two signs differ, the local factor is the
domain of lines of one prescribed sign in a Hermitian plane of
signature \((1,1)\).

For \(t\in\mathcal T\), the two compact-place pieces of the
positive Lagrangian are
\[
 V_\tau\otimes W_\tau^-,
 \qquad
 V_\tau^*\otimes(W_\tau^+)^*.
\]
In the second expression \((W_\tau^+)^*\) is embedded as
\(\operatorname{ann}(W_\tau^-)\subset W_\tau^*\); the
orthogonal decomposition identifies it with the dual of
\(W_\tau^+\).  The active pieces remain, exactly,
\[
 L\otimes W_\mu,\qquad Q^*\otimes W_\mu^*.
\]
They depend only on \(W_\mu\), not on any decomposition of it.
These prescriptions give positive Lagrangians for all
\(t\in\mathcal T\).  The same projection as in
\eqref{eq:theta-top-target}, with a local smooth choice of
constant-line trivialization, defines a smooth form of type
\((2,0)\), denoted \(\Xi_W(\phi;t)\), even when \(t\) is not
induced by rational lines.  We do not claim it is holomorphic
for such an intermediate parameter.

\begin{proposition}\label{prop:theta-top}
The rank-two construction has the following properties.
\begin{enumerate}[label=(\roman*)]
\item At period data induced by a rational orthogonal
  decomposition \(W=W_1\perp W_2\), the form \(\Xi_W(\phi)\)
  is holomorphic.  On a sufficiently deep common principal
  quotient it is a finite sum, with constant coefficients, of
  forms \(\Theta_1(\phi_{1,r})\wedge\Theta_2(\phi_{2,r})\).
\item There is an input which is a product in adapted rational
  coordinates and for which \(\Xi_W(\phi)\) is nonzero.
\item Fix \(W\), the total finite input \(\phi\), and a
  torsion-free principal quotient \(X\) whose finite level fixes
  \(\phi\) under \(\omega_{W,f}\).  With locally smooth
  constant-line trivializations, the map
  \[
    t\longmapsto\Xi_W(\phi;t)
  \]
  is continuous in the \(C^\infty\) topology of differential
  forms on this fixed compact quotient.  The same finite level
  works throughout the parameter neighborhood.
\end{enumerate}
\end{proposition}

\begin{proof}
For (i), the rational orthogonal decomposition gives a rational
symplectic direct sum
\(\mathbb V_W=\mathbb V_{W_1}\perp\mathbb V_{W_2}\).
Choose rational coordinates adapted to it.  By
Remark~\ref{rem:theta-adapted-coordinates}, the total reference
input becomes another finite Schwartz input in the adapted
coordinates.  Such an input is a finite sum
\[
 \sum_r\phi_{1,r}\otimes\phi_{2,r}.
\]
Indeed, a locally constant compactly supported function on the
product of two finite-adelic vector spaces is invariant under a
sufficiently small product compact open subgroup.  Only finitely
many of its cosets meet the support, and each coset is a product
coset.

For a product input the completed theta series multiply, since
both the sum and its completion split in the adapted coordinates.
If \(T_j^{(i)}\) denotes the degree-\(j\) coefficient for the
\(i\)-th line, the coefficient of total degree two is
\begin{equation}\label{eq:theta-quadratic-product}
 T_2^{(1)}T_0^{(2)}
       +T_1^{(1)}T_1^{(2)}
       +T_0^{(1)}T_2^{(2)}.
\end{equation}
The coefficient of the mixed product is \(1\), by the
\(1/k!\) normalization in \eqref{eq:theta-completion}.
The two unmixed terms have both multiplicity factors in the
same one-dimensional \((W_i)_\mu^*\); they vanish under
\(\operatorname{alt}_2\).  The mixed term becomes the wedge
of the two projected first tensors.  With compatible
constant-line trivializations it is exactly
\(\Theta_1(\phi_{1,r})\wedge\Theta_2(\phi_{2,r})\);
other trivializations multiply it by a fixed nonzero scalar.
Rational changes of coordinates are already accounted for by
\eqref{eq:theta-taylor-law}.  Moreover, two identifications
with the differential bundle used here have a ratio invariant
under the transitive infinite special-unitary action on \(D\),
so that ratio is constant.  This proves the wedge expansion
and its holomorphy.  We have proved holomorphy only after the
double-alternating projection: the whole completed second
Taylor tensor need not be holomorphic.  The finitely many
individual inputs descend on a further common principal level
by Lemma~\ref{lem:unitary-splitting}.

For (ii), fix a point of \(D\).  Proposition~\ref{prop:theta-one}
provides an input for \(W_1\) with a nonzero cotangent value and
an input for \(W_2\) whose value is outside the line it spans.
Their wedge is nonzero at that point.  The product of these
inputs in adapted coordinates has a nonzero double-alternating
second tensor by the calculation above.  Transporting this
product back to the reference coordinates gives a total finite
input with \(\Xi_W(\phi)\ne0\).

For (iii), no rational coordinate change is made as \(t\)
varies.  The rational symplectic space, its rational \(G\)-action,
the finite representation, and the total input all remain
fixed.  The changed compact-place Lagrangians remain invariant
under \(\SU(3)\), since all changes are in the multiplicity
factor.  Their determinant characters on \(\SU(3)\) are
trivial.  The active pieces have not changed, so their
determinant character remains \(\ell^4\).  The infinite
comparison in Lemma~\ref{lem:unitary-splitting} therefore gives
the same half-determinant action \(L^2\) at every \(t\).
Consequently the same subgroup fixing \(\phi\) gives
equivariance and descent at every parameter.  There is no
finite translation, Hecke choice, or change of level in this
variation.

It remains to check uniformity on the quotient.  The constant
determinant factors in \eqref{eq:theta-halfdet} vary smoothly
with the compact-place subspaces.  On a small parameter
neighborhood choose smooth nonzero frames and their square
roots, agreeing with the chosen frames at the center.
Cover the fixed compact \(X\) by finitely many charts with
precompact lifts to \(D\).  On the closure of each lifted chart
times a sufficiently small compact parameter neighborhood, the
period map is smooth, its image is a compact subset of the
Siegel domain, and every fixed order of its derivatives is
bounded.  The imaginary parts thus have a common positive
lower bound.  The rational support of the fixed \(\phi\)
continues to lie in one fixed finite union of lattice cosets.
The Gaussian estimates following
\eqref{eq:theta-completion} now give uniform convergence of
every fixed order of derivatives in the chart coordinates and
the parameters.  Projection to \eqref{eq:theta-top-target}
and the chosen smooth trivializations preserve these bounds.
The unchanged equivariance glues these convergent local
expressions on \(X\).  This is continuity in the \(C^\infty\)
topology.  Compactness of \(X\) is used here; no estimate at a
cusp or boundary is needed.
\end{proof}

\subsection{Approximation while the total theta input stays fixed}\label{subsec:theta-approximation}

We now apply the last proposition to the sign switch.  The
approximation occurs in \(\mathcal T\), before the rational
coordinate changes used to factor either endpoint.

\begin{corollary}\label{cor:theta-switch-continuity}
Suppose \(d_1+d_2=d_3+d_4\), with \(d_i(\mu)=1\) for
all \(i\).  Start with \(W=W_1\perp W_2\) as above, let
\(t_0\) be its compact-place data, and choose a total input
\(\phi\) for which
\(\xi=\Xi_W(\phi;t_0)\) is a nonzero holomorphic top form
on a sufficiently deep compact principal quotient \(X\).
There are rational orthogonal decompositions
\(W=W_3\perp W_4\) with line signs \(d_3,d_4\) whose
total compact-place positive Lagrangians are arbitrarily close
to those at \(t_0\).  The total active positive Lagrangian is
exactly unchanged.

For a sufficiently close rational choice \(t'\), the form
\(\xi'=\Xi_W(\phi;t')\) is holomorphic on the same quotient
\(X\), with the same total finite input, and
\begin{equation}\label{eq:theta-nonzero-pairing}
 \int_X\xi\wedge\overline{\xi'}\ne0.
\end{equation}
Both endpoints admit the wedge expansions of
Proposition~\ref{prop:theta-top} on a further common principal
cover, and their pairing remains nonzero there.
\end{corollary}

\begin{proof}
At a compact place \(\tau\ne\mu\), equality of the two sign
sums says exactly that the multisets
\(\{d_1(\tau),d_2(\tau)\}\) and
\(\{d_3(\tau),d_4(\tau)\}\) coincide.  Reorder the two
original local lines so that the first has the sign required
by \(d_3(\tau)\), and choose a nonzero vector on it.
After fixing an \(E\)-basis of \(W\), weak approximation in
\(E\), coordinate by coordinate at these finitely many
complex places, gives a single \(w\in W(E)\) arbitrarily
close to all the chosen local vectors.  Nondegeneracy and
the signs of their Hermitian norms are open conditions.
For a sufficiently close approximation,
\[
 W_3=Ew,\qquad W_4=W_3^\perp
\]
is therefore a rational orthogonal decomposition.  Its first
line has the prescribed compact-place signs for \(d_3\);
orthogonal inertia gives the complementary signs for \(d_4\).
At \(\mu\) the form on \(W_\mu\) is definite, so both new
lines automatically have the same required sign there.

The total compact-place Lagrangians converge to the original
ones.  Reordering the original lines only relabels their total
positive summands.  For a nondegenerate local vector the
orthogonal projection onto its line is
\[
 z\longmapsto h_W(z,w)h_W(w,w)^{-1}w.
\]
It varies continuously with \(w\), so both the line and its
orthogonal complement converge to the chosen original ones.
The selected signed subspaces, and hence the total positive
Lagrangians, converge.  At the active pair the total pieces
are always \(L\otimes W_\mu\) and \(Q^*\otimes W_\mu^*\),
so they are unchanged, not merely close.  The approximating
rational lines may have small active norms or large arithmetic
denominators.  No bound on these is needed: the active total
period data are fixed, and the continuity assertion concerns
the total tensor in the original reference coordinates.

All the hypotheses of Proposition~\ref{prop:theta-top}(iii)
now hold with the same rational \(W\), the same \(G\), the
same total finite \(\phi\), and the same compact \(X\).
Thus
\[
 \Xi_W(\phi;t')\longrightarrow \Xi_W(\phi;t_0)
\]
smoothly on \(X\).  The limit \(\xi\) is a nonzero
holomorphic top form.  In local complex coordinates
\(\xi=f\,dz_1\wedge dz_2\), the form
\(\xi\wedge\bar\xi\) is a positive multiple of
\(|f|^2\,dx_1\wedge dy_1\wedge dx_2\wedge dy_2\).
Consequently
\(\int_X\xi\wedge\bar\xi>0\).  Continuity of the integral
gives \eqref{eq:theta-nonzero-pairing} for all sufficiently
close rational choices.  At each such endpoint holomorphy
follows from Proposition~\ref{prop:theta-top}(i), not from a
holomorphy assertion along the intermediate parameter path.

Factor the two endpoints only now.  There are finitely many
line inputs in their two wedge expansions.  Choose a principal
subgroup contained both in the level of \(X\) and in their
finite stabilizers.  It gives a finite unramified cover
\(\pi:X'\to X\) on which all individual factors descend.
The covering formula is
\begin{equation}\label{eq:theta-cover-pairing}
 \int_{X'}\pi^*\xi\wedge\overline{\pi^*\xi'}
       =\deg(\pi)\int_X\xi\wedge\overline{\xi'}\ne0.
\end{equation}
Both quotients are compact and torsion-free, so there is no
cusp term or ramification correction.  This proves the last
assertion.
\end{proof}

Expansion of the two pulled-back endpoint forms on \(X'\) in
\eqref{eq:theta-cover-pairing} gives
\eqref{eq:theta-mixed-period-goal}; in particular at least one
pair of individual wedge terms has nonzero mixed period.
This is the analytic nonvanishing needed for
Proposition~\ref{prop:surface-switch}.
To apply that criterion
we must still express the relevant one-form classes as complex
linear combinations of images of the selected CM lines under
rational Hodge maps, on a common projective principal cover.
That arithmetic step is separate from the approximation just
proved.
